\documentclass[11pt,a4paper]{article}
\usepackage[utf8]{inputenc}
\usepackage[T1]{fontenc}
\usepackage{amsmath,amssymb,amsthm}
\usepackage{amsfonts}
\usepackage{mathtools}
\usepackage{geometry}
\usepackage{hyperref}
\usepackage{enumitem}
\usepackage{booktabs}
\usepackage{array}
\usepackage{mathrsfs}
\usepackage{xcolor}
\usepackage{microtype}
\usepackage{fancyhdr}
\usepackage{lastpage}

\geometry{a4paper, top=2.5cm, bottom=2.5cm, left=2.5cm, right=2.5cm}

\pagestyle{fancy}
\fancyhf{}
\renewcommand{\headrulewidth}{0pt}
\fancyfoot[C]{\thepage\ / \pageref{LastPage}}

\DeclareMathOperator{\End}{End}
\DeclareMathOperator{\Sp}{Sp}
\DeclareMathOperator{\ad}{ad}
\DeclareMathOperator{\Ker}{Ker}
% \Im is already defined in amssymb
\DeclareMathOperator{\Grass}{Grass}

\begin{document}

\begin{flushright}
1987
\end{flushright}

\vspace{1cm}

\begin{center}
\Large\textbf{Letter to Millson}
\end{center}

\vspace{0.5cm}

\noindent Dear Millson,

\vspace{0.5cm}

Nori brought me a copy of your paper with Goldman on deformations of flat bundles. It seems to me that it contains a number of mistakes, but also that, when corrected, it gives the strong results you ask for $p\,8\,l.\,10\,\grave{a}\,-7$. I begin with the latter.

\vspace{0.3cm}

The philosophy, which I had not realized before reading your paper, seems the following: in characteristic $0$, a deformation problem is controlled by a differential graded Lie algebra (d.g.l.a.), with quasi-isomorphic DG Lie algebras giving the same deformation theory. If the DG Lie algebra controlling a problem is ``formal'', i.e.\ quasi-isomorphic to $\oplus H^i$, then the versal (formal) deformation space is that of $\oplus H^i$, i.e.\ is the completion at $0$ of the subscheme of $H^1$ defined by the equations $[u,u]=0$.

\vspace{0.3cm}

If one believes in this, what you do becomes transparent:

\vspace{0.2cm}

\noindent\textbf{Deformation problem:} $G$ is a given reductive group, and one wants to deform a given $G$-torsor $P$ --- i.e.\ a $G(\mathbb{C})$-torsor, or, equivalently, a $G(\mathbb{C})$-torsor, with integrable connection $\nabla$.

\vspace{0.2cm}

\noindent\textbf{Controlling DG algebra:} (one is over $\mathbb{C}$):
\[
\Omega^{**}(\operatorname{Lie} G^P)
\]

If $P$ can be reduced to a maximal compact $K \subset G$ (i.e.\ to a compact form of $G$), Hodge theory shows that the controlling

\newpage

\noindent algebra is ``formal'': one has quasi-isomorphisms of DG Lie algebras
\begin{equation}\label{eq:1}
\Omega^{**}(\operatorname{Lie} G^P) \longleftrightarrow \Ker(d') \longrightarrow \frac{\Ker(d')}{\Im(d')}
\end{equation}
and $\Ker(d')/\operatorname{Im}(d')$ has $d=0$. To check the above, it is easiest to use the principle of the 2 types in the form that the double complex $\Omega^{**}(\operatorname{Lie} G^P)$ is (forgetting $[,]$) a sum of

\vspace{0.2cm}
\begin{tabbing}
\hspace{1cm}\= a)\hspace{0.3cm}\= a double complex with $d'=d''=0$ \hspace{1cm}\= (``harmonic part'')\\[0.2cm]
\> b)\> squares of isomorphisms: \> (to shewline)
\end{tabbing}
\vspace{-0.3cm}
\[
\begin{array}{c}
* \xrightarrow{\;d'\;} * \\[0.1cm]
{\scriptstyle d''}\Big\updownarrow \quad\; \Big\updownarrow{\scriptstyle d''} \\[0.1cm]
* \xrightarrow[\;d'\;]{}
\end{array}
\]

As far as formal deformations are concerned, this concludes the story. It is not clear to me how to use Artin in the analytic case [sol.\ of anal.\ eqn.] to get the same for analytic-versal. The algebraic case is OK, but Artin is more powerful there.

\vspace{0.3cm}

Here is how to give a meaning to the philosophy (and check it):

\vspace{0.2cm}

\noindent\textbf{Construction:} How to attach to a DG Lie algebra $L^*$ (over a chosen field $k$) a fibred category over the category of $\Sp(A)$, $A$ local $k$-algebras of finite dimension with residue field $k$.

\vspace{0.2cm}

\noindent\textbf{Object over $\Sp(A)$:} $[k = A/m]$
\[
\omega \in L^1 \otimes m, \quad \text{such that} \quad d\omega + \tfrac{1}{2}[\omega,\omega] = 0
\]

\vspace{0.2cm}

\noindent\textbf{Map:} The Lie algebra $L^0$ gives rise to a formal group $\widehat{\underline{L}}^0$; the $A$-points of this formal group are the $L^0 \otimes m$, viewed as a nilpotent Lie algebra, and equipped with Campbell--Hausdorff group law.

\noindent Notation: $\widehat{\underline{L}}^0(A)$.

\newpage

The formal group $\widehat{\underline{L}}^0$ acts on $L^1$ (thanks to $[,]\colon L^0 \otimes L^1 \to L^1$), hence $\widehat{\underline{L}}^0(A)$ on $L^1 \otimes A$. One has also a map
\[
\text{``}g^{-1}dg\text{''}\colon \widehat{\underline{L}}^0 \longrightarrow L^1 \quad\colon\quad \exp(X) \longmapsto \frac{1-\exp(-\ad X)}{\ad X}\,[dX]
\]
i.e.
\[
\widehat{\underline{L}}^0(A) \longrightarrow L^1 \otimes m
\]

One now defines a map from $\omega \in L^1 \otimes m$ to $\omega' \in L^1 \otimes m$ as being $g \in \widehat{\underline{L}}^0(A)$ such that
\begin{equation}\label{eq:2}
\omega' = g\,dg^{-1} + g(\omega)
\end{equation}
One has to check that $d\omega' + \frac{1}{2}[\omega',\omega'] = g\bigl(d\omega + \frac{1}{2}[\omega,\omega]\bigr)$ when $\omega$ and $\omega'$ are related by (2).

\vspace{0.3cm}

\noindent\textbf{To be checked:} for a square zero ideal $I$ in $A$, and an object $\omega$ on $A/I$, there is an obstruction in $H^2(L \otimes I)$ to extend $\omega$ to $A$; if the obstruction vanishes, isomorphism classes of extensions form a principal homogeneous space under $H^1(L \otimes I)$; automorphisms of any extension (trivial on $A/I$) are $H^0(L^* \otimes I)$. This is functorial in $L^*$, hence the fact that quasi-isomorphisms induce equivalences of fibred categories.

\vspace{0.3cm}

\noindent\textbf{Application here:} Let $X$ be an analytic variety and $(V,\nabla)$ be a (holomorphic) vector bundle with connection. The deformation problem is
\[
A \longmapsto \text{category of vector bundles with connection in the $X$-direction,}
\]
on $X \otimes \Sp(A)$, deforming $(V,\nabla)$.

It is given by the previous construction applied to $\Omega^{**}(\End V)$.

\newpage

\noindent\textbf{Rmq:} In your application, because you have the explicit quasi-isomorphism (1), you have an explicit equivalence of fibred categories between
\begin{itemize}[leftmargin=2cm]
\item[$*$] deformations of $(V,\nabla)$ as above
\item[$*$] category with objects over $\Sp(A)$: $\Sp(A) \longrightarrow$ subscheme of $H^1_L$ with $[u,u]=0$

and arrows: $g \in (\underline{H}^0_L)^\wedge(A)$ from $u$ to $g(u)$
\end{itemize}

\vspace{0.3cm}

The functor $T$ is such: a cohomology class $u \in H^1_L \otimes m$, with $\varepsilon_u\colon \Sp(A) \to \text{scheme } H^1_L$, can be uniquely written as $\omega \in L^1 \otimes m$ with $d'\omega = d''\omega = 0$. This $\omega$ can be corrected by a $d''\beta$ to have zero curvature, if $[u,u]=0$. The corrected $\tilde{\omega}$ is unique mod $\exp(\Ker d' \otimes m)$.

\vspace{0.3cm}

\noindent\textbf{Rmq:} (1) is a filtered quasi-isomorphism, for the Hodge filtration. This should translate into saying that the pair of spaces
\[
\bigl(\text{deformations of }(V,\nabla)\bigr) \supset \bigl(\text{deformations of }\nabla\text{ on a fixed }V\bigr)
\]
is isomorphic to the pair
\[
\bigl(H^1 \quad\supset\quad F^0 H^1\bigr) \cap \Grass\,[u,u]=0.
\]

\newpage

\section*{II. Comments on your paper}

\begin{itemize}[leftmargin=0.5cm]
\item[$p\,2$] after cor.\ 3: no [take $G = \mathrm{U}_1$]
\item[$p\,6$] l.\,18 holomorphic? : I guess you want to start with $\eta$ such that $d'\eta = d''\eta = 0$
\item[$p\,7$] l.\,4 : why closed $\Rightarrow$ (2,0)-component closed?
\end{itemize}

\vspace{0.5cm}

\section*{III. Generalization}

Hodge theory is available for variations of Hodge structures (complex variations, on compact K\"ahler manifolds). The useful bigrading of $\Omega^{**}(\End V)$ is obtained by mixing that of $\Omega^*$ and that of $\End V$. It is useful now that (1) is a filtered quasi-isomorphism. One gets:

\vspace{0.2cm}

\begin{itemize}[leftmargin=1cm]
\item[a)] The same result as before on deformations of a representation of $\pi_1$ given by a complex variation.
\item[b)] Looking at the deformation theory controlled by $F^0(\cdots)$: same result for the versal moduli of a filtered bundle with connection coming from a complex variation. One insists that the deformed connection must continue to satisfy ``transversality'':
\[
\nabla F^i \subset \Omega^1 \otimes F^{i-1}.
\]
\end{itemize}

\vspace{1.5cm}

\begin{flushright}
Yours sincerely,

\vspace{0.8cm}
P.\,Deligne

\vspace{0.2cm}
P.\,DELIGNE
\end{flushright}

\end{document}
